\beforeabstract
\prefacesection{Abstract}

Wearable fall detectors are trained almost exclusively on falls simulated by healthy volunteers,
yet they are meant for people whose movement, and whose falls, look different. We study the most
extreme version of that mismatch: a deep fall detector trained on the SisFall dataset (38 healthy
young and elderly subjects) and applied, without any patient label, to freezing of gait (FoG) in the
Daphnet dataset (10 patients with Parkinson's disease). The comparison is controlled: one accelerometer at the same anatomical site in both domains, rotated
into a common vertical/antero-posterior/medio-lateral frame identified from the data; both domains in
$g$ at \SI{64}{\hertz}; fall labels centred on the detected impact; and identical seeds, mini-batches,
loss, optimiser and tuning budget for a CNN-1D, an LSTM, a GRU and a patch-based Transformer.
In-domain, the four models are almost indistinguishable (AUROC 0.988--0.999, recall 0.96--0.98,
0.7--1.5 false alarms per hour on unseen subjects). Under zero-shot transfer, AUROC drops to
0.49--0.63; only the Transformer's subject-level confidence interval excludes chance. At their validated
thresholds the detectors stay silent on FoG (recall below 1\%) while false alarms on walking rise 15- to
18-fold. A training-free freeze index ranks FoG better (AUROC 0.66) than every network. Label-free
re-normalisation of the target inputs gives no consistent gain. In every learned embedding, a linear
probe still separates healthy from Parkinsonian movement (AUROC $\geq$ 0.96), and the fall direction
is nearly orthogonal to the FoG direction. Fine-tuning from fall weights is worse than training from
scratch in six of eight comparisons. The transfer gap is therefore dominated by concept shift rather
than covariate shift, and in-domain accuracy gives no indication of transferability.

\textbf{Keywords:} domain shift, fall detection, freezing of gait, Parkinson's disease, wearable
sensors, Transformer, recurrent neural networks.

\prefacesection{Acknowledgements}

This work grew out of a group project carried out at ENSA Marrakech with Rokhaya Gueye and Agrhib
Azar, whom I thank for their collaboration and discussions.
I thank Pr.\ Mustapha Ameur for supervising the project. The SisFall dataset was made available by the
SISTEMIC group of the Universidad de Antioquia, and the Daphnet dataset by the DAPHNet project and the
Tel Aviv Sourasky Medical Center.

\afterpreface \afterabstract
